\documentclass[aps,prl,article,twocolumn,preprintnumbers,amsmath,amssymb,superscriptaddress]{revtex4-2}

\usepackage{graphicx}
\usepackage{dcolumn}
\usepackage{bm}
\usepackage{amsmath}
\usepackage{amssymb}
\usepackage{mathrsfs}
\usepackage{epigraph}
\usepackage{gensymb}
\usepackage{lmodern}
\usepackage{tipa}
\usepackage{bbold}
\usepackage{esint}
\usepackage{mathdots}
\usepackage{xcolor}
\usepackage{appendix}
\usepackage{natbib}
\usepackage{multirow}
\usepackage{float}
\usepackage{physics}
\usepackage{comment}
\usepackage[colorlinks=true, linkcolor=blue, urlcolor=blue, citecolor=blue]{hyperref}

\begin{document}

\appendix


\section{Derivation of the Electronic Free-Energy Model}\label{app:derivationFreeEnergy}

In this section, we derive the effective electronic free-energy model
used in the main text from a microscopic interacting electron model.
The low-energy electrons are described by
\begin{equation}
H
=
H_0
+
H_{\mathrm{int}},
\end{equation}
where
\begin{equation}
H_0
=
\sum_{\mathbf k,n,s}
\epsilon_{n,\mathbf k,s}
d^\dagger_{n,\mathbf k,s}
d_{n,\mathbf k,s}
\end{equation}
is the noninteracting Hamiltonian, where $d_{n,\mathbf k,s}$ is the
annihilation operator on a band basis and $s=\uparrow,\downarrow$ labels
spin.

To capture the low-energy density-wave-like ordering tendency at finite
momentum, we consider an effective momentum-conserving interaction in
the density channel,
\begin{equation}
\label{eq:siInteraction}
H_{\mathrm{int}}
=
-
\frac12
\sum_{\mathbf q}
V_{\mathbf q}
\rho_{-\mathbf q}
\rho_{\mathbf q},
\end{equation}
where $V_{\mathbf q}$ is the effective interaction strength.
The density operator projected to the band basis is
\begin{equation}
\rho_{\mathbf q}
=
\sum_{\mathbf k,m,n,s}
g_{mn}(\mathbf k,\mathbf q)
d^\dagger_{m,\mathbf k+\mathbf q,s}
d_{n,\mathbf k,s},
\end{equation}
with the band form factor
\begin{equation}
g_{mn}(\mathbf k,\mathbf q)
=
\left\langle
u_{m,\mathbf k+\mathbf q}
\middle|
u_{n,\mathbf k}
\right\rangle .
\end{equation}
The form factor encodes the wave-function overlap between the two
electronic states connected by momentum transfer $\mathbf q$.

To derive the effective free energy, we introduce a Hubbard--Stratonovich
field $\psi_{\mathbf q}$ in the same density channel under static mean-field,
\begin{equation}
e^{-\beta H_{\mathrm{int}}}
=
\prod_{\mathbf q}
\int
\mathcal D[\psi_{\mathbf q}]
e^{-S[d,d^\dagger,\psi_{\mathbf q}]},
\end{equation}
where the auxiliary field is 
\begin{equation}
\label{eq:siHSField}
S[d,d^\dagger,\psi_{\mathbf q}]
=
\frac{
|\psi_{\mathbf q}|^2
}{
2V_{\mathbf q}
}
+
\psi_{-\mathbf q}\rho_{\mathbf q},
\end{equation}
where $\psi_{\mathbf q}$ represents the finite-momentum electronic order
parameter.  This form can be understood by integrating out the HS field,
which recovers the original density-density interaction in
Eq.~\eqref{eq:siInteraction}.  Equivalently, at the saddle point one
obtains $\psi_{\mathbf q}\propto -V_{\mathbf q}\rho_{\mathbf q}$,
showing that a nonzero $\psi_{\mathbf q}$ is the same density-wave order
described by $\rho_{\mathbf q}$.  The overall sign is conventional and
can be absorbed into the phase of $\psi_{\mathbf q}$.

After the HS transformation, the electron action is quadratic in the
fermion fields.  The field $\psi_{\mathbf q}$ acts as a static scattering
potential that mixes electronic states with momenta $\mathbf k$ and
$\mathbf k+\mathbf q$.  Integrating out the electronic degrees of freedom
therefore yields
\begin{equation}
e^{-\beta \mathcal{F}_{\mathrm{ele}}[\psi]}
=
\int
\mathcal D[d,d^\dagger]
e^{-S[d,d^\dagger,\psi]} ,
\end{equation}
or, equivalently,
\begin{equation}
\mathcal F_{\mathrm{ele}}[\psi]
=
\sum_{\mathbf q}
\frac{|\psi_{\mathbf q}|^2}{2V_{\mathbf q}}
-
T\,\mathrm{Tr}
\ln
\left(
G_0^{-1}
+
\Sigma_\psi
\right).
\end{equation}
Here $G_0$ is the Green's function of $H_0$, $\Sigma_\psi$ is the
scattering vertex generated by $\psi_{\mathbf q}$, and $\mathrm{Tr}$
denotes a trace over momentum, Matsubara frequency, band, and spin indices.
Expanding the logarithm gives
\begin{equation}\label{siLogExpansion}
-
T\mathrm{Tr}
\ln
\left(
1
+
G_0\Sigma_\psi
\right)
=
-
T
\sum_{\ell=1}^{\infty}
\frac{(-1)^{\ell+1}}{\ell}
\mathrm{Tr}
\left(
G_0\Sigma_\psi
\right)^\ell .
\end{equation}
The linear term vanishes for a finite ordering wavevector because one
insertion of $\Sigma_\psi$ changes the electron momentum by
$\pm\mathbf q$ and cannot close the fermion loop.  The first nonzero term
from the trace-log expansion is therefore quadratic and represents the
bubble.  Keeping only this term, we obtain
\begin{equation}
\label{eq:siQuadraticTrace}
\mathcal F_{\mathrm{ele}}^{(2)}
=
\sum_{\mathbf q}
\frac{|\psi_{\mathbf q}|^2}{2V_{\mathbf q}}
+
\frac{T}{2}
\mathrm{Tr}
\left(
G_0\Sigma_\psi
G_0\Sigma_\psi
\right).
\end{equation}
For a static density-wave order, $\Sigma_\psi$ contains the scattering
vertex
\begin{equation}
\left(\Sigma_\psi\right)_{m,\mathbf k+\mathbf q;n,\mathbf k}
=
\psi_{-\mathbf q}
g_{mn}(\mathbf k,\mathbf q),
\end{equation}
together with its Hermitian-conjugate process.  In spin space the vertex
is diagonal,
\begin{equation}
\left(\Sigma_\psi\right)_{m,\mathbf k+\mathbf q,s;n,\mathbf k,s'}
=
\psi_{-\mathbf q}
g_{mn}(\mathbf k,\mathbf q)
\delta_{ss'} .
\end{equation}

To identify the trace term, we first recall the density response function
in the normal state,
\begin{equation}
\chi(\mathbf q,\tau)
=
-
\left\langle
T_\tau
\rho_{\mathbf q}(\tau)
\rho_{-\mathbf q}(0)
\right\rangle_{0},
\end{equation}
where the expectation value is evaluated using $H_0$.  Substituting the
projected density operator into this expression and applying Wick's
theorem gives the Matsubara-frequency bubble
\begin{align}
\chi(\mathbf q,i\Omega_l)
=&
-
T
\sum_{\omega_n,\mathbf k,m,n,s}
\left|
g_{mn}(\mathbf k,\mathbf q)
\right|^2
\nonumber\\
&\times
G_{0,m,s}(i\omega_n+i\Omega_l,\mathbf k+\mathbf q)
G_{0,n,s}(i\omega_n,\mathbf k).
\end{align}
The spin index is the same on the two Green's functions, reflecting that
only spin-conserving particle-hole excitations contribute to the charge
susceptibility.  Equivalently, $\chi=\chi_\uparrow+\chi_\downarrow$.
Substituting the same vertex into the trace in
Eq.~\eqref{eq:siQuadraticTrace} gives precisely this density bubble.
Therefore, in the static limit relevant for the
mean-field free energy,
\begin{align}\label{eq:siBubbleTrace}
\frac{T}{2}
\mathrm{Tr}
\left(
G_0\Sigma_\psi
G_0\Sigma_\psi
\right)
=&
-
\frac{1}{2}
\sum_{\mathbf q}
\chi(\mathbf q;T,n)
|\psi_{\mathbf q}|^2 ,
\end{align}
where $\chi(\mathbf q;T,n)\equiv \chi(\mathbf q,i\Omega_l=0)$ after the
usual analytic continuation.  The minus sign is the fermion-loop sign in
the bubble definition.  Equivalently, the quadratic part of the free
energy can be written as
\begin{equation}
\mathcal F_{\mathrm{ele}}^{(2)}
=
\sum_{\mathbf q}
\frac{1}{2}
\left(\frac{1}{V_{\mathbf q}}
-
\chi(\mathbf q;T,n)
\right)
|\psi_{\mathbf q}|^2 .
\end{equation}
The factor $1/2$ appears because the full momentum summation counts the
pair $\mathbf q$ and $-\mathbf q$ twice, while
$\psi_{-\mathbf q}=\psi_{\mathbf q}^{*}$ for a real density modulation.
In the main text we keep one independent ordering wavevector, so this
overall factor is absorbed into the convention for the Landau functional.

Including the next nonvanishing term in the trace-log expansion gives the
quartic contribution.  In the main text we keep only its symmetry-allowed
form, $b_{\mathbf q}|\psi_{\mathbf q}|^4$, and take $b_{\mathbf q}>0$ so
that the free energy is stable.  The resulting Ginzburg--Landau free
energy is
\begin{equation}
\mathcal F_{\mathrm{ele}}
=
\sum_{\mathbf q} \left[
a(\mathbf q;T,n)
|\psi_{\mathbf q}|^2
+
b_{\mathbf q}
|\psi_{\mathbf q}|^4
+
\mathcal{O}(|\psi_{\mathbf q}|^6) \right],
\end{equation}
where $a(\mathbf q;T,n) = 1/V_{\vb*{q}} - \chi_{\vb*{q}}(T,n)$.

We now evaluate the static susceptibility entering this coefficient.
The noninteracting Green's function is
\begin{equation}
G_{0,n,s}(i\omega_n,\mathbf k)
=
\frac{1}{i\omega_n-\epsilon_{n,\mathbf k,s}}
\end{equation}
and the Matsubara frequency sum in the bubble is
\begin{align}
&T
\sum_{\omega_n}
G_{0,m,s}(i\omega_n+i\Omega_l,\mathbf k+\mathbf q)
G_{0,n,s}(i\omega_n,\mathbf k)
\nonumber\\
&\hspace{1.0cm}
=
\frac{
f(\epsilon_{n,\mathbf k,s})
-
f(\epsilon_{m,\mathbf k+\mathbf q,s})
}{
i\Omega_l
+
\epsilon_{n,\mathbf k,s}
-
\epsilon_{m,\mathbf k+\mathbf q,s}
}.
\end{align}
Taking the static limit after analytic continuation,
$i\Omega_l\rightarrow \omega+i\eta$ and then $\omega\rightarrow 0$,
therefore leads to the bubble form used in the main text,
\begin{equation}
\chi(\mathbf q;T,n)
=
\sum_{\mathbf k,m,n,s}
\left|
g_{mn}(\mathbf k,\mathbf q)
\right|^2
\frac{
f(\epsilon_{n,\mathbf k,s})
-
f(\epsilon_{m,\mathbf k+\mathbf q,s})
}{
\epsilon_{m,\mathbf k+\mathbf q,s}
-
\epsilon_{n,\mathbf k,s}
+
i\eta
}.
\end{equation}

\section{Model Parameters for the Lattice Potential}\label{app:latticePotential}

For the hysteresis calculations, we use a single real lattice coordinate
$X$ for the lattice distortion amplitude.  The lattice free energy is
modeled by a local sixth-order Landau potential,
\begin{equation}
\label{eq:siLatticePotential}
\mathcal F_{\mathrm{lat}}(X)
=
\frac{1}{2}a_2 X^2
+
\frac{1}{4!}b_2 X^4
+
\frac{1}{6!}c_2 X^6 .
\end{equation}
The sixth-order term stabilizes the potential when $b_2<0$.  In the
code, this form is implemented in \texttt{free\_energy.m}, where the
electronic order parameter is \texttt{psi1} and the lattice coordinate
is \texttt{psi2}.  The full free energy minimized in the hysteresis
calculation is
\begin{align}
\label{eq:siNumericalFreeEnergy}
\mathcal F(\psi,X)
=&
\frac{1}{2}a_1\psi^2
+
\frac{1}{4!}b_1\psi^4
\nonumber\\
&+
\frac{1}{2}a_2 X^2
+
\frac{1}{4!}b_2 X^4
+
\frac{1}{6!}c_2 X^6
+
\lambda \psi X .
\end{align}
Here $\psi$ is the electronic density-wave order parameter and $X$ is
the lattice order parameter.  The bilinear coupling $\lambda\psi X$ is
the single-mode version of the electron-phonon-induced coupling derived
in Sec.~\ref{app:ephElectronicCorrection}.

For the hysteresis calculation, the electronic quadratic coefficient is
obtained from the calculated static susceptibility as
\begin{equation}
\label{eq:siA1Numerical}
a_1(T,n)
=
s_\chi
\left[
V^{-1}
-
|\chi(\mathbf q;T,n)|
\right],
\end{equation}

The lattice coefficients are taken to be functions of temperature and
doping,
\begin{align}
\label{eq:siLatticeCoeff}
a_2&=a_2(T),\\
b_2&=b_2(n),\\
c_2&=\mathrm{const.}
\end{align}
Here $T$ is measured in Kelvin and $n$ is the doping coordinate used in
the numerical scan.
With these conventions, the lattice transition criterion tracked in the
phase-diagram code is
\begin{equation}
\Delta_{\mathrm{lat}}
=
\frac{5}{6}b_2^2
-
a_2c_2
=0 ,
\end{equation}
which locates the onset of nonzero-$X$ local minima for the sixth-order
lattice potential.

The numerical values used in the hysteresis calculation are summarized
in Table~\ref{tab:siHystParams}.
\begin{table}[h]
\caption{
Parameters used in the hysteresis free-energy calculation.
}
\label{tab:siHystParams}
\begin{ruledtabular}
\begin{tabular}{cc}
parameter & value \\
\hline
$V^{-1}$ & $0.069$ \\
$s_\chi$ & $300$ \\
$b_1$ & $1$ \\
$\lambda$ & $2$ \\
$a_2(T)$ & $0.2\,T$ \\
$b_2(n)$ & $-\sqrt{\max[0,0.24(10-2n^2)]}$ \\
$c_2$ & $1$ \\
\end{tabular}
\end{ruledtabular}
\end{table}

\section{Electronic Correction from Electron-Phonon Coupling}
\label{app:ephElectronicCorrection}

We now include the coupling between the electronic density and a static
lattice displacement $X_{\mathbf q}$.  The total Hamiltonian becomes
\begin{equation}
H
=
H_0
+
H_{\mathrm{int}}
+
H_{\mathrm{e\text{-}ph}},
\end{equation}
where $H_{\mathrm{int}}$ is still the density-density interaction in
Eq.~\eqref{eq:siInteraction}.  The electron-phonon coupling is taken to
be in the same density channel,
\begin{equation}
H_{\mathrm{e\text{-}ph}}
=
\sum_{\mathbf q}
-g_{\mathbf q}X_{-\mathbf q}\rho_{\mathbf q},
\end{equation}
where $g_{\mathbf q}$ is the electron-phonon coupling constant.  In this
section $X_{\mathbf q}$ is treated as a given lattice coordinate and is
not integrated out.

Since $H_{\mathrm{e\text{-}ph}}$ couples linearly to the same density
operator as the HS field, it can be included in the HS transformation as
an external source.  The interaction and source part is decoupled as
\begin{equation}
e^{-\beta (H_{\mathrm{int}}+H_{\mathrm{e\text{-}ph}})}
=
\prod_{\mathbf q}
\int
\mathcal D[\psi_{\mathbf q}]
e^{-S[d,d^\dagger,\psi_{\mathbf q},X_{\mathbf q}]},
\end{equation}
where the shifted action is
\begin{equation}
S[d,d^\dagger,\psi_{\mathbf q},X_{\mathbf q}]
=
\frac{
\left|
\psi_{\mathbf q}
+
g_{\mathbf q}X_{\mathbf q}
\right|^2
}{
2V_{\mathbf q}
}
+
\psi_{-\mathbf q}\rho_{\mathbf q}.
\end{equation}
This equation is the same HS identity as before, but shifted by the
external source $g_{\mathbf q}X_{\mathbf q}$.  Integrating out the HS
field recovers both the attractive density-density interaction and the
linear electron-phonon coupling.
Using the same expansion in Eq.~\eqref{siLogExpansion} and keeping the
bubble term, we obtain
\begin{equation}
\mathcal F_{\mathrm{ele}}
=
\frac{
\left|
\psi_{\mathbf q}
+
g_{\mathbf q}X_{\mathbf q}
\right|^2
}{
2V_{\mathbf q}
}
+
\frac{T}{2}
\mathrm{Tr}
\left(
G_0\Sigma_\psi
G_0\Sigma_\psi
\right).
\end{equation}

The important point is that the phonon displacement does not enter the
fermionic bubble directly at this order.  After the HS transformation, the
electrons still see the scattering vertex generated by
$\psi_{\mathbf q}$,
\begin{equation}
\left(\Sigma_\psi\right)_{m,\mathbf k+\mathbf q,s;n,\mathbf k,s'}
=
\psi_{-\mathbf q}
g_{mn}(\mathbf k,\mathbf q)
\delta_{ss'} .
\end{equation}
Therefore the trace-log expansion gives the same susceptibility correction
as in Eq.~\eqref{eq:siBubbleTrace}
\begin{align}
\frac{T}{2}
\mathrm{Tr}
\left(
G_0\Sigma_\psi
G_0\Sigma_\psi
\right)
=&
-
\frac{1}{2}
\sum_{\mathbf q}
\chi(\mathbf q;T,n)
|\psi_{\mathbf q}|^2 .
\end{align}
the correction to the Ginzburg--Landau free energy can be written as
\begin{align}
\Delta \mathcal F_{\mathrm{ele}}^{(ph)}[\psi,X]
=
&-
\sum_{\mathbf q}
\frac{1}{2V_{\mathbf q}}
g_{\mathbf q}X_{\mathbf q}\psi_{-\mathbf q}
+
\sum_{\mathbf q}
\frac{|g_{\mathbf q}X_{\mathbf q}|^2}{2V_{\mathbf q}},
\end{align}
where the quadratic coefficient of the lattice order can be absorbed into the lattice free energy model


\section{First-Principles Simulations}\label{app:DFT}



\bibliographystyle{apsrev4-2}
\bibliography{refs}

\end{document}
